\BourbakiExerciseGroup

\begin{enumerate}
\def\labelenumi{\arabic{enumi})}
\item
  设 E 是域 K 的一个扩张，x 与 y 是 K{[}X{]} 中同一不可约多项式在 E 中的两个不同的根。证明扩张 \(\mathbf{K}(x)\) 与 \(\mathbf{K}(y)\) 在 K 上不是线性不相交的（利用定理 2）。* 给出一个例子，其中 K = Q 且 \(\mathbf{K}(x) \cap \mathbf{K}(y) = \mathbf{K}\) （参见 V，第 161 页，习题 2）。*
\item
  设 \((\mathbf{E}_i)_{i\in I}\) 是域 \(\mathbf{K}\) 的一族扩张，都包含于 \(\mathbf{L}\) （ \(\mathbf{K}\) 的一个扩张）中。若 \(\mathbf{F}_i\) 是 \(\mathbf{K}\) 在 \(\mathbf{E}_i\) 中的相对代数闭包，证明 \(\mathbf{K}\) 在 \(\mathbf{E} = \bigcap_{i\in I}\mathbf{E}_i\) 中的相对代数闭包是 \(\mathbf{F} = \bigcap_{i\in I}\mathbf{F}_i\) .
\item
  证明域 K 的每个代数扩张 E 都与 \(K \times N\) 的一个子集等势（考虑把 E 的每个元素对应到它在 K 上的极小多项式的映射）。特别地，有限域的每个代数扩张都是可数的，而无限域 K 的每个代数扩张都与 K 等势。* 由此推出，在实数域 R 中存在素子域 Q 上的超越数，并且所有这些数组成的集合具有连续统的势。*
\item
  设 F 是域 E 的一个超越扩域。证明
\end{enumerate}

\[
[ \mathrm{F}: \mathrm{E} ] \geqslant \operatorname * {S u p} (\operatorname{Card} (\mathrm{E}), \operatorname{Card} (\mathbf {N}))
\]

若扩张 F 可由一族可数的元素生成，则等号成立。（利用 IV，第 89 页，习题 7。）

\begin{enumerate}
\def\labelenumi{\arabic{enumi})}
\setcounter{enumi}{4}
\tightlist
\item
  设 A 是一个交换代数，S 是 A 的一个子集，它作为 K-代数生成 A。假设 \(\operatorname{Card}(S) < \operatorname{Card}(K)\) 。证明：若 M 是 A 的一个极大理想，则域 A/M 是 K 的一个代数扩张。（当 \(\operatorname{Card}(K) \leqslant \operatorname{Card}(N)\) 时利用 V，第 174 页，习题 15，当
\end{enumerate}

\[
\operatorname{Card} (\mathbf {K}) > \operatorname{Card} (\mathbf {N}).
\]

时利用习题 4。）* 由此推出，\(\mathbf{A} / \mathfrak{M} = \mathbf{K}\) （当 \(\mathbf{K}\) 是代数闭域时）。*

* 6) 设 C 是一个不可数的代数闭域，且 \(\mathbf{P} = \mathbf{C}[(\mathbf{X}_i)_{i \in I}]\) 是 C 上的一个多项式代数。给定 P 的元素的两个族 \((f_\alpha)_{\alpha \in A}, (g_\beta)_{\beta \in B}\) ，假设：a) B 与 I 都是可数的。

\begin{enumerate}
\def\labelenumi{\alph{enumi})}
\setcounter{enumi}{1}
\tightlist
\item
  对于每个有限子集 \(\mathbf{A}'\) （ \(\mathbf{A}\) 的）与每个有限子集 \(\mathbf{B}'\) （ \(\mathbf{B}\) 的），存在 \((x_i)_{i\in I} \in C^I\) 使得
\end{enumerate}

\[
f _ {\alpha} (x _ {i}) = 0 \quad \text { and } \quad g _ {\beta} (x _ {i}) \neq 0 \quad \text { for   all } \quad \alpha \in \mathbf {A} ^ {\prime}, \beta \in \mathbf {B} ^ {\prime}.
\]

证明此时存在 \((x_{i})_{i\in I}\in \mathbf{C}^{I}\) 使得

\[
f _ {\alpha} (x _ {i}) = 0 \quad \text { and } \quad g _ {\beta} (x _ {i}) \neq 0 \quad \text { for   all } \quad \alpha \in A, \beta \in B.
\]

（设 \((\mathbf{Y}_{\beta})_{\beta \in \mathbb{B}}\) 是以 \(\mathbf{B}\) 为指标集的未定元族，并设 \(\mathbf{Q}\) 为多项式代数 \(\mathbf{C}[(\mathbf{X}_i)_{i\in I}, (\mathbf{Y}_{\beta})_{\beta \in \mathbb{B}}]\) 。设 \(\Lambda\) 为 \(\mathbf{Q}\) 关于由诸 \(f_{\alpha}\) 与诸 \(1 - g_{\alpha}\mathbf{Y}_{\alpha}\) 生成的理想所成的商。证明 \(\Lambda \neq 0\) （借助 \(a)\) 与 \(b)\) ），并对 \(\Lambda\) 的一个极大理想应用习题 5。）*

\begin{enumerate}
\def\labelenumi{\arabic{enumi})}
\setcounter{enumi}{6}
\tightlist
\item
  设 \(L\) 是一个域， \(A\) 是 \(L\) 的一个子环，且 \(K \subset L\) 是 \(A\) 的分式域.
\end{enumerate}

\begin{enumerate}
\def\labelenumi{\alph{enumi})}
\item
  证明：若 L 是有限生成的 A-模（II，第 206 页），则必有 A = K（考虑 K 在 L 中的补子空间（视 L 为 K 上的向量空间），证明 K 是有限生成的 A-模，从而必具有形式 \(s^{-1}A\) ，其中 \(s \in A\) ；最后证明 s 在 A 中必可逆）。
\item
  证明：若存在属于 L 的有限个元素 \(x_{j}\) ( \(1 \leqslant j \leqslant n\) )，它们在 K 上是代数的，使得 \(L = A[x_{1}, \ldots, x_{n}]\) ，则存在 A 中的元素 \(b \neq 0\) 使得 \(K = A[b^{-1}]\) （证明存在属于 A 的 \(b \neq 0\) ，使得 L 是有限生成的 \(A[b^{-1}]\) -模）。证明：要么 b 可逆（此时 K = A），要么 A 的每个非零理想都包含一个主理想 \(Ab^{m}\) 。反之亦然。域 k 上的形式幂级数环 \(k[[X]]\) 就是此类环的一个例子.
\end{enumerate}

* 8) 若 K 是域，证明 K 在域 K(X)（K 上单变元有理分式域）中是代数闭的（利用 K{[}X{]} 是主理想整环这一事实）。*

\begin{enumerate}
\def\labelenumi{\arabic{enumi})}
\setcounter{enumi}{8}
\tightlist
\item
  设 \(\mathbf{K}\) 是一个域， \(\mathbf{E}\) 是 \(\mathbf{K}\) 的一个扩张，且 \(\mathbf{S}\) 是 \(\mathbf{E} - \mathbf{K}\) 的一个子集.
\end{enumerate}

\begin{enumerate}
\def\labelenumi{\alph{enumi})}
\item
  证明由扩张 \(\mathbf{L} \subset \mathbf{E}\) （ \(\mathbf{K}\) 的、满足 \(\mathbf{L} \cap \mathbf{S} = \varnothing\) 的）组成的集合按包含关系排序具有一个极大元 \(\mathbf{M}\) .
\item
  证明：若 S 是有限的，则 E 是 M 的代数扩张。
\end{enumerate}

* 10) 设 K 是一个域，P 是环 K{[}X, Y{]} 中两个未定元的多项式。假设存在多项式 H ∈ K{[}X{]} 以及两个多项式 Q, R ∈ K{[}X, Y{]} 使得 H(X) P(X, Y) = Q(X, Y) R(X, Y)，并且 Q(X, Y) 的系数（作为 Y 的多项式）是 K{[}X{]} 中的多项式，且没有非常数的公因子。证明：R(X, Y) 的所有系数（作为 Y 的多项式）都能被 H(X) 整除。（将 H, P, Q, R 视为域 K(Y) 上关于 X 的多项式，并将它们在 Ω{[}X{]} 中分解为一次因式，其中 Ω 是 K(Y) 的一个代数闭包；注意，若 H(X) 与 Q(X, Y) 有一个一次公因子，则该因子将整除 Q(X, Y) 的所有系数（作为 Y 的多项式），并利用事实（IV，第 12 页）：若 K' 是 K 的一个扩域，则 K'{[}X{]} 中两个 K{[}X{]} 的多项式的最大公因式属于 K{[}X{]}。）*

\begin{enumerate}
\def\labelenumi{\arabic{enumi})}
\setcounter{enumi}{10}
\tightlist
\item
  设 \(\mathbf{E}\) 是域 \(\mathbf{K}, x \in \mathbf{E}\) 是 \(\mathbf{K}\) 上的一个超越元，从而 \(\mathbf{K}(x)\) 同构于 \(\mathbf{K}(\mathbf{T})\) （即 \(\mathbf{K}\) 上关于一个未定元的有理分式域）。于是每个 \(y \in \mathbf{K}(x)\) 都可写成 \(y = g(x) / h(x)\) ，其中 \(g\) 与 \(h\) 是 \(\mathbf{K}[\mathbf{X}]\) 的两个互素的多项式（IV，第 12 页），在相差 \(\mathbf{K}\) 中一个非零公因子的意义下确定；\(g\) 与 \(h\) 的次数中的较大者称为 \(y\) 关于 \(x\) 的高度.
\end{enumerate}

\begin{enumerate}
\def\labelenumi{\alph{enumi})}
\tightlist
\item
  证明在多项式环 \(\mathbf{K}(y)[\mathbf{X}]\) 中，多项式 \(g(\mathbf{X}) - yh(\mathbf{X})\) 当 \(y \notin \mathbf{K}\) 时是不可约的。由此推出，若 \(y\) 的高度为 \(n > 0\) （关于 \(x\) 的），则 \(\mathbf{K}(x)\) 是一个 \(n\) 次的、关于 \(\mathbf{K}(y)\) 的代数扩张.
\end{enumerate}

设 \(\mathrm{P}(\mathbf{X}) = \sum_{i=0}^{n} y_i \mathbf{X}^i\) 为 \(x\) 在 \(\mathrm{K}(y)\) 上的极小多项式。证明对于

\(0 \leqslant i \leqslant n\) ，我们或者有 \(y_i \in \mathbf{K}\) ，或者有 \(\mathbf{K}(y_i) = \mathbf{K}(y)\) .

\begin{enumerate}
\def\labelenumi{\alph{enumi})}
\setcounter{enumi}{1}
\item
  由 a) 推出，每个 \(y \in \mathbf{K}(x)\) 若满足 \(\mathbf{K}(y) = \mathbf{K}(x)\) ，则具有形式 \((ax + b)/(cx + d)\) ，其中 a、b、c、d 是 K 中满足 \(ad - bc \neq 0\) 的元素；反之亦然。求出 \(\mathbf{K}(x)\) 的全部 K-自同构.
\item
  证明：若 \(y \in \mathbf{K}(x)\) 的高度为 \(n\) （关于 \(x\) 的），且 \(z \in \mathbf{K}(y)\) 的高度为 \(m\) （关于 \(y\) 的），则 \(z\) 的高度为 \(mn\) （关于 \(x\) 的）.
\end{enumerate}

* \(d\) ) 设 F 是 K 的一个扩域，使得 \(\mathbf{K} \subset \mathbf{F} \subset \mathbf{K}(x)\) 且 \(\mathbf{F} \neq \mathbf{K}\) ；设 y 是 F 中的一个元素，其高度 \(m\) （关于 \(x\) ）取到可能的最小值；证明 \(\mathbf{F} = \mathbf{K}(y)\) （“吕罗特定理”）。（设 P 为 \(x\) 在 F{[}T{]} 中的极小多项式；证明 \(\mathrm{P(T)} = \mathrm{Q(T,x)}/\mathrm{R}(x)\) ，其中 Q 是 K{[}T,X{]} 中的一个多项式，它关于 X 的次数 \(\geqslant m\) ，且不被 K{[}X{]} 的任何非常数多项式整除，而 R 是 K{[}X{]} 的一个多项式。若 \(y = g(x)/h(x)\) ，其中 \(g\) 与 \(h\) 在 K{[}X{]} 中互素，注意由 \(a\) )，多项式 \(g(T)h(X) - g(X)h(T)\) 不被 K{[}T{]} 或 K{[}X{]} 的任何非常数多项式整除；由此必然推出

\[
g (T) h (X) - g (X) h (T) = c. Q (T, X), \quad \text { where } c \in K,
\]

利用习题10。）*

* 12) a) 采用习题 11 中的记号与假设，证明：K 的扩张 F 满足 \(\mathbf{K} \subset \mathbf{F} \subset \mathbf{K}(x)\) ，要使 \(\mathbf{F} \cap \mathbf{K}[x] \neq \mathbf{K}\) ，其充要条件是 \(\mathbf{F} = \mathbf{K}(y)\) ，其中 \(y = g(x)\) ，\(g \in \mathbf{K}[\mathrm{T}]\) 为非常数多项式；此时我们有 \(\mathbf{K}(y) \cap \mathbf{K}[x] = \mathbf{K}[y]\) 。 （写 \(\mathbf{F} = \mathbf{K}(y)\) ，其中 \(y = g(x)/h(x)\) , \(g\) 与 \(h\) 是 \(\mathbf{K}[\mathrm{T}]\) 的互素多项式。如果存在一个非常数多项式 \(\mathbf{P} \in \mathbf{K}[\mathrm{T}]\) 使得 \(\mathbf{P}(x) \in \mathbf{F}\) ，我们可以写出 \(\mathbf{P}(x) = \mathbf{Q}(y)/\mathbf{R}(y)\) ，其中 \(\mathbf{Q}\) 与 \(\mathbf{R}\) 是 \(\mathbf{K}[\mathrm{T}]\) 的两个互素的首一多项式，并非两者皆为常数。分两种情况讨论，视 \(\mathbf{R} = 1\) 或 \(\mathbf{R}\) 为非常数而定；在第二种情形下，将 \(\mathbf{Q}\) 与 \(\mathbf{R}\) 在 \(\Omega\) （ \(\mathbf{K}\) 的一个代数闭包）中分解为一次因式，并注意到，如果 \(a, b\) 是 \(\Omega\) 的不同元素，则 \(g - ah\) 与 \(g - bh\) 在 \(\Omega[\mathrm{T}]\) 中互素。）

\begin{enumerate}
\def\labelenumi{\alph{enumi})}
\setcounter{enumi}{1}
\item
  如果 \(\mathbf{F} = \mathbf{K}(y)\) ，其中 \(y = g(x)\) ，这里 \(g \in \mathbf{K}[T]\) 是非常数的，则我们可以假设 \(g(T) = Tg_1(T)\) ，其中 \(g_1 \in \mathbf{K}[T]\) 。证明如果 \(\mathbf{K}(xg(x)) \cap \mathbf{K}[y] \neq \mathbf{K}\) ，我们必然有 \(g_1(T) = aT^n\) ，其中 \(a \in \mathbf{K}\) 且 \(n > 0\) 。（注意，根据假设及 \(a\) ），存在两个非常数多项式 \(P, Q\) （在 \(K[T]\) 中）使得 \(P(xg(x)) = Q(g(x)) \notin K\) ；由此推出存在两个元素 \(a, b\) （在 \(K\) 中）以及一个整数 \(m\) 使得 \(g(T)\) 整除 \(a - bT^m\) 。） c) 设 \(y, z\) 为属于 \(K[x]\) 的两个元素，使得 \(K(y) \cap K(z) \neq K\) ；证明此时 \(K[y] \cap K[z] \neq K\) .
\item
  求 \(y \in \mathbf{K}(x)\) 使得 \(\mathbf{K}(y) = \mathbf{K}(x)\) 但是 \(\mathbf{K}[y] \cap \mathbf{K}[x] = \mathbf{K}\) . *
\end{enumerate}

* 13) 设 K 为一个域， \(P(T, X) = a_0(T)X^n + a_1(T)X^{n-1} + \cdots + a_n(T)\) 为 K{[}T, X{]} 中的一个多项式；假设： \(1^\circ\) 多项式 \(a_0\) 不被 T 整除； \(2^\circ\) 多项式 \(a_1, \ldots, a_n\) 中的每一个都被 T 整除； \(3^\circ\) 多项式 \(a_n\) 不被 \(T^2\) 整除。证明在这些条件下，P 作为 K(T){[}X{]} 中的多项式是不可约的（用反证法，利用习题 10）。*

* 14：采用习题 11 中的记号，证明如果我们取 \(g(X) = X\) 且 \(h(X) = X^n + X + 1 \quad (n \geqslant 2)\) ，则多项式 \((g(X)h(x) - h(X)g(x)) / (X - x)\) （作为 \(K(x)[X]\) 的元素）是不可约的（利用习题 13）。*

* 15）给出一个域 K 的例子，它包含两个子域 \(\mathbf{K}_1\) , \(\mathbf{K}_2\) ，使得 K 在 \(\mathbf{K}_1\) 上是代数的，在 \(\mathbf{K}_2\) 上也是代数的，但在 \(\mathbf{K}_1 \cap \mathbf{K}_2\) 上不是代数的。 （取 \(\mathbf{K} = \mathbf{Q}(T)\) , \(\mathbf{K}_1 = \mathbf{Q}(T^2)\) ，以及 \(\mathbf{K}_2 = \mathbf{Q}(T^2 - T)\) 。验证 \([\mathbf{K} : \mathbf{K}_1] = 2\) , \([\mathbf{K} : \mathbf{K}_2] = 2\) 且 \(\mathbf{K}_1 \cap \mathbf{K}_2 = \mathbf{Q}\) 。 为此，注意一个元素 \(f(T)\) （ \(\mathbf{K}_1 \cap \mathbf{K}_2\) 中的）满足条件 \(f(-T) = f(T)\) 与 \(f(1 - T) = f(T)\) .)

